\documentclass[a4paper, noanswers, addpoints, 12pt]{exam}
\usepackage[T1]{fontenc}
\usepackage[utf8]{inputenc}
\usepackage{lmodern}
\usepackage{tikz}
\usetikzlibrary{positioning,shapes.multipart,arrows.meta,shadows}
\usepackage[french]{babel}
\usepackage{xcolor}
\usepackage{listings}
\usepackage{graphicx}
\usepackage[most]{tcolorbox}

\definecolor{codegreen}{rgb}{0,0.6,0}
\definecolor{codegray}{rgb}{0.5,0.5,0.5}
\definecolor{codepurple}{rgb}{0.58,0,0.82}
\definecolor{backcolour}{rgb}{0.95,0.95,0.92}
 
\lstdefinestyle{mystyle}{
    backgroundcolor=\color{backcolour},   
    commentstyle=\color{codegreen},
    keywordstyle=\color{magenta},
    numberstyle=\tiny\color{codegray},
    stringstyle=\color{codepurple},
    basicstyle=\footnotesize,
    breakatwhitespace=false,         
    breaklines=true,                 
    captionpos=b,                    
    keepspaces=true,        
    numbers=left,                    
    numbersep=5pt,                  
    showspaces=false,                
    showstringspaces=false,
    showtabs=false,                 
    tabsize=2
}
\lstset{style=mystyle, language=Python}

\newcommand{\squote}{\textquotesingle}

\title{CC1 Term NSI - Rattrapage : Pile, File et Rappels de 1er}
\author{Babash - NSI Tani Malandi et Johanne Bakalara}
\date{Septembre 2025, Durée: 45 minutes}


\usepackage[a4paper,top=2.3cm,bottom=2.8cm,left=2cm,right=2cm,headsep=0.6cm,footskip=1.2cm]{geometry}
\lstset{extendedchars=true,literate={è}{{\`e}}1 {é}{{\'e}}1 {ê}{{\^e}}1}
\begin{document}
\runningfooter{}{\footnotesize Babash - NSI Tani Malandi et Johanne Bakalara \textperiodcentered{} CC BY-SA 4.0 \textperiodcentered{} page \thepage}{}
\firstpagefooter{}{\footnotesize Babash - NSI Tani Malandi et Johanne Bakalara \textperiodcentered{} CC BY-SA 4.0 \textperiodcentered{} page \thepage}{}

\maketitle
\thispagestyle{headandfoot}
\begin{center}\footnotesize\itshape D'après l'exercice 1 du sujet 22-NSIJ1G11 (baccalauréat 2022), le sujet 22-NSIJ2ME1 (baccalauréat 2022, Métropole, jour 2) et l'exercice 1 du sujet 21-NSIJ2G11 (baccalauréat 2021).\end{center}

\pagestyle{headandfoot}
\thispagestyle{headandfoot}
\firstpageheader{\makebox[\textwidth]{Nom : \dotfill\ Pr\'enom : \dotfill\ Classe : \makebox[3cm]{\dotfill}}}{}{}


\section*{Consignes:}
\begin{center} 
\fbox{\parbox{\linewidth}{
\begin{center}
Le devoir est à rédiger sur une feuille libre.\\
La clarté de la rédaction et les explications font partie intégrante de la notation.\\
Une réponse non rédigée sans explication sera considérée fausse.\\
Prenez le temps d'encadrer les réponses.
\end{center}
}}
\end{center}

\vsword{Note}
\hsword{Note}

\section*{Barème :}
\begin{center}
  \par\bigskip
  \gradetable[h]
\end{center}

\vspace{0.5cm}

\section{Exercice I : Dictionnaires et Listes (4 points)}
\textit{Thème abordé : structures de données : dictionnaire et liste}\\

On dispose de la liste \textit{jours} et du dictionnaire \textit{mois} suivants :
\begin{lstlisting}[language=Python]
jours = ["dimanche", "lundi", "mardi", "mercredi",
         "jeudi", "vendredi", "samedi"]
mois = {1 :("janvier",31), 2 :("février",28), 3 :("mars",31),
        4 :("avril",30)}
\end{lstlisting}

\vspace{0.5cm}

\begin{questions} 
\question Manipulation de liste et dictionnaire
\begin{parts}
\part[1] Écrire le code Python pour obtenir "jeudi" à partir de la liste \textit{jours}.
\part[1] Écrire le code Python permettant d'obtenir le \textbf{nom} du mois numéro \textit{2} à partir du dictionnaire \textit{mois}.
\end{parts}

\vspace{0.25cm}
\question Pour chaque affirmation, choisir la \textbf{seule} bonne réponse.
\begin{parts}
    \part[1] Quel est le principal avantage d'utiliser un \textbf{dictionnaire} plutôt qu'une \textbf{liste} ?
    \begin{itemize}
        \item[A.] La liste est plus rapide.
        \item[B.] Le dictionnaire permet un accès direct par clé (nom/numéro) sans parcourir les éléments.
        \item[C.] Le dictionnaire peut contenir des éléments en double, contrairement à la liste.
    \end{itemize}

    \part[1] Que renvoie l'instruction Python \texttt{jours[7]} ?
    \begin{itemize}
        \item[A.] \texttt{'samedi'}
        \item[B.] \texttt{'dimanche'}
        \item[C.] Une erreur de type \texttt{IndexError}
    \end{itemize}

    \part[1] Une \textbf{liste} Python est un objet \textbf{mutable}. Qu'est-ce que cela signifie ?
    \begin{itemize}
        \item[A.] Ses éléments ne peuvent être que des nombres.
        \item[B.] Ses éléments ne peuvent pas être modifiés après sa création.
        \item[C.] Ses éléments ou sa taille peuvent être modifiés après sa création.
    \end{itemize}

    \part[1] Si on exécute l'instruction \texttt{mois[5] = ('mai', 31)}, le dictionnaire \textit{mois} compte alors :
    \begin{itemize}
        \item[A.] 4 couples clé/valeur.
        \item[B.] 5 couples clé/valeur.
        \item[C.] Une erreur car on ne peut pas ajouter d'éléments à un dictionnaire après sa création.
    \end{itemize}
\end{parts}

\vspace{0.25cm}

\question[2] Fonction avec Dictionnaire
Écrire une fonction \textit{nb\_jours\_mois(mois, numero\_mois)} qui prend en paramètre le dictionnaire \textit{mois} et le numéro d'un mois, et qui renvoie le \textbf{nombre de jours} de ce mois.
Par exemple : \textit{nb\_jours\_mois(mois, 3)} doit renvoyer \textit{31}.
\vspace{0.25cm}

\question[4] Code à trou : Changement de mois
Recopier et compléter le code à trou suivant pour une fonction \textit{changement\_mois(numero\_mois, annee)} qui prend le numéro du mois actuel \textit{numero\_mois} et l'année \textit{annee}, et renvoie un tuple correspondant au mois suivant (\textit{numero\_mois\_suivant, annee\_suivante}).

Par exemple:
\begin{itemize}
    \item \texttt{changement\_mois(1, 2025)} renvoie \texttt{(2, 2025)}
    \item \texttt{changement\_mois(12, 1900)} renvoie \texttt{(1, 1901)}
\end{itemize}

\begin{lstlisting}[language=Python]
def changement_mois (numero_mois, annee):
    if numero_mois == ...:
        return (..., annee + ...)
    else:
        return (numero_mois + ..., annee)
\end{lstlisting}

\newpage

\section{Exercice II : Pile et File (6 points)}

\question Question de cours : LIFO/FIFO
\begin{parts}
    \part[2] Que signifie l'acronyme \textbf{FIFO} ? Caractérise-t-il une pile ou une file ?
    \part[1] Donner un exemple simple d'une \textbf{file} dans la vie courante.
\end{parts}

\vspace{0.5cm}

\question[2] Réduction de la pile par exécution de code
On considère la pile $P$ (du haut vers le bas) : $\mathbf{5, 4, 5, 4, 2, 1}$.

La fonction \textit{parcourir\_pile\_en\_reduisant(p)} réalise un parcours de la pile du haut vers le bas en appliquant la règle de réduction.

Nous considérons les fonctions (dont l'interface de la pile est supposée connue) :

\begin{lstlisting}[language=Python]
def reduire_triplet_au_sommet(p):
    # Règle: Supprime l'élément central (b) si le haut (a) et le bas (c)
    # des 3 premiers éléments sont de même parité.
    if taille(p) >= 3:
        a = depiler(p)
        b = depiler(p)
        c = sommet(p)
        if a % 2 != c % 2:
            # Pas de réduction: on remet b et a
            empiler(p, b)
            empiler(p, a)
        else:
            # Réduction: on remet juste a (b est supprimé)
            empiler(p, a) 

def parcourir_pile_en_reduisant(p):
    q = creer_pile_vide()
    while taille(p) >= 3:
        reduire_triplet_au_sommet(p)
        e = depiler(p)
        empiler(q, e)
    while not est_vide(p): # On vide la fin de p dans q
        empiler(q, depiler(p))
        
    while not est_vide(q): # On remet tous les éléments de q dans p (dans le bon ordre)
        empiler(p, depiler(q))
    return p
\end{lstlisting}

Après l'appel de l'instruction \texttt{parcourir\_pile\_en\_reduisant(P)}, quel est le contenu de la pile $P$ (du haut vers le bas) ?

\vspace{0.5cm}

\question[3] Fonction \textit{est\_gagnant}
Une partie est \og gagnante\fg\ lorsque la pile finale est réduite à \textbf{deux éléments exactement}.

Écrire une fonction \textit{est\_gagnant(pile\_finale)} qui prend en paramètre une pile \textit{pile\_finale} et renvoie \textit{True} si la partie est gagnante et \textit{False} sinon.\\
Indice : On pourra utiliser la fonction \textit{taille(p)} pour connaitre la taille de la pile \textit{p}.

\end{questions}

\end{document}